\providecommand{\toplevelprefix}{..}  %
\documentclass[../book-main_es.tex]{subfiles}
\begin{document}
\chapter*{Acrónimos}\label{ch:acronyms}

%\begin{longtable}{|l|l|l|}
%\begin{longtable}{@{} p{\dimexpr 0.15\linewidth - 2\tabcolsep\relax} @{\hspace{1em}} p{\dimexpr 0.55\linewidth - 2\tabcolsep\relax} @{\hspace{1em}} p{\dimexpr 0.55\linewidth - 2\tabcolsep\relax} @{}}
\begin{longtable}{@{} p{\dimexpr 0.12\linewidth - 2\tabcolsep\relax} @{\hspace{1em}} p{\dimexpr 0.55\linewidth - 2\tabcolsep\relax} @{\hspace{1em}} p{\dimexpr 0.55\linewidth - 2\tabcolsep\relax}}
\caption{Acrónimos} \label{tab:acronyms} \\
\hline
Sigla & Expansión & Traducción \\
\hline
\endfirsthead
\multicolumn{3}{l}{\small (sigue)} \\
\hline
Sigla & Expansión & Traducción \\
\hline
\endhead
AD & Automatic Differentiation & Diferenciación automática \\
ADM & Ablated Diffusion Model & Modelo de difusión con ablación \\
AI & Artificial Intelligence & Inteligencia artificial \\
ANN & Artificial Neural Network & Red neuronal artificial \\
AoT & Attention only Transformer & Transformer de solo atención \\
AP & Average Precision  & Precisión promedio \\
ARI & Adjusted Rand Index & Índice Rand ajustado \\
ASLDM & Aligned Shape Latent Diffusion Model & Modelo latente de difusión con formas alineadas \\
BP & Back Propagation & Retropropagación \\
BPE & Byte Pair Encoding & Codificación de pares de bytes \\
CE & Cross-entropy & Entropía cruzada \\
CFG & Classifier-free Guidance & Guía sin clasificador \\
CFM & Conditional Flow Matching & Matching de flujos condicionales \\
CG & Classifier Guidance & Guía por clasificador \\
CIFAR & Canadian Institute for Advanced Research & Instituto Canadiense de Investigación Avanzada \\
CISE & NSF Directorate for Computer and Information Science and Engineering & Directorio NSF para la Ingeniería y Ciencia de la Información y la Computación \\
CLIP & Contrastive Language–image Pre-training & Pre-entrenamiento contrastivo de lenguaje-imagen \\
CNN & Convolutional Neural Network & Red neuronal convolucional \\
CRATE & Coding Rate Transformer & Transformer de tasa de codificación \\
CS & Computer Science & Ciencias de la computación  \\
CTRL & Conditional Transformer Language Model & Modelo de lenguaje Transformer condicional \\
CV & Computer Vision & Visión artificial \\
DAPS & Decoupled Annealing Posterior Sampling & Muestreo a posteriori mediante recocido desacoplado \\
DDIM & Denoising Diffusion Implicit Model & Modelo de difusión implícita para la eliminación de ruido \\
DDPM & Denoising Diffusion Probabilistic Model & Modelo probabilístico de difusión de eliminación de ruido \\
DINO & Distillation with No labels & Destilación sin etiquetas \\
DiT & Diffusion Image Transformer & Transformer de imagen de difusión \\
DL & Dictionary Learning & Aprendizaje de diccionarios \\
DMS & NSF Data Management and Sharing Plan & Plan NSF de Gestión y Compartición de Datos \\
DNA & Deoxyribonucleic acid & Ácido desoxirribonucleico \\
DNN & Deep Neural Network & Red neuronal profunda \\
DPS & Diiffusion Posterior Sampling & Muestreo a posteriori de difusión \\
ELBO & Evidence Lower Bound & Cota inferior de evidencia \\
EM & Earth Mover’s Distance (EMD) & Distancia del transportista de tierra, también Distancia de Wasserstein \\
EMA & Exponential Moving Average & Media móvil exponencial \\
FID & Fréchet Inception Distance & Distancia de inicio de Fréchet \\
FM & Flow Matching & Coincidencias de flujo \\
FPS & Farthest Point Sampling & Muestreo del punto más lejano \\
GAN & Generative Adversarial Network & Red generativa adversativa \\
GCN & Graph Convolutional Network & Red gráfica convolucional \\
GD & Gradient Descent & Descenso de gradiente \\
GDA & Gradient Descent-Ascent & Descenso-ascenso de gradiente \\
GELU & Gaussian Error Linear Unit & Unidad lineal de error gaussiano \\
GMM & Gaussian Mixture Model & Modelo de mezcla gaussiana \\
GPCA & Generalized Principal Component Analysis & Análisis  de componentes principales generalizado \\
GPU & Graphics Processing Unit & Unidad de procesamiento gráfico \\
HaMeR & Hand Mesh Recovery & Recuperación de malla de mano \\
HMD & Head Mounted Device/Display & Dispositivo/Visor montado en la cabeza \\
IB & Information Bottleneck & Cuello de botella de la información \\
ICA & Independent Component Analysis & Análisis  de componentes independientes \\
IoU & Intersection over Union & Intersección sobre unión \\
ISTA & Iterative Shrinkage-Thresholding Algorithm & Algoritmo iterativo de reducción y umbralización \\
JEPA & Joint Embedding Predictive Architecture & Arquitectura de predicción con embeddings conjuntos \\
JS & Jensen-Shannon & de Jensen-Shannon \\
KL & Kullback–Leibler & de Kullback–Leibler \\
LASSO & Least Absolute Shrinkage and Selection Operator & Operador de selección y compresión absoluta mínima \\
LDM & Latent Diffusion Model & Modelo de difusión latente \\
LDR & Linear Discriminative Representations & Representaciones lineales discriminantes \\
LISTA & Learned Iterative Shrinkage-Thresholding Algorithm & Algoritmo iterativo aprendido de reducción y umbralización \\
LLM & Large Language Model & Modelo de lenguaje de gran tamaño \\
LTI & Linear Time-Invariant & Lineal e invariante en el tiempo \\
LUCA & Last Universal Common Ancestor & Último antepasado común universal \\
MAE & Masked Autoencoder & Autocodificador enmascarado \\
MAP & Maximum a posteriori  & Maximo a posteriori  \\
MCL & Minimizing Coding Length & Minimización de la longitud del código \\
MCR² & Maximal Coding Rate Reduction & Reducción máxima de la tasa de codificación \\
MDL & Minimum Description Length & Longitud mínima de descripción \\
MHSA & Multi-Head Self-Attention & Autoatención multicabezal \\
ML & Machine Learning & Aprendizaje automático \\
MLP & Multi Layer Perceptron & Perceptrón multicapa \\
MML & Minimum Message Length & Longitud mínima de mensaje \\
MMSE & Minimum Mean-Square Error & Error cuadrático medio mínimo \\
MNIST & Modified National Institute of Standards and Technology & Instituto Nacional Modificado de Estándares y Tecnología \\
MoE & Mixture of Experts & Mezcla de expertos \\
MPJPE & Mean Per-Joint Position Error & Error medio de posición por articulación \\
MRI & Magnetic Resonance Imaging & Imaginería de resonancia magnética \\
MSE & Mean-Square Error & Error cuadrático medio \\
MSP & Matching, Stretching, and Projection & Coincidencia, estiramiento y proyección \\
MSSA & Multihead Subspace Self-Attention  & Autoatención multicabezal por subespacios \\
NAS & Neural Architecture Search & Búsqueda de arquitectura neuronal \\
NLP  & Natural Language Processing & Procesamiento del lenguaje natural \\
NLPCA & Nonlinear Principal Component Analysis & Análisis no lineal de componentes principales \\
NN & Neural Network & Red neuronal  \\
ODE & Ordinary Differential Equation & Ecuación diferencial ordinaria \\
ODL & Overcomplete Dictionary Learning & Aprendizaje de diccionarios sobrecompletos \\
PCA & Principal Component Analysis & Análisis de componentes principales \\
PDE & Partial Differential Equation & Ecuación diferencial parcial \\
PDF & Probability Density Function & Función de densidad de probabilidad \\
PPCA & Probabilistic Principal Component Analysis & Análisis probabilístico de componentes principales \\
PSD & Positive Semi-Definite & Semidefinido positivo \\
RAE & Representation Autoencoder & Autocodificador de representaciones \\
RAE & Recurrent Autoencoder & Autocodificador recurrente \\
RD & Rate-Distortion & Tasa-distorsión \\
ReLU & Rectified Linear Unit & Unidad lineal rectificada \\
RGB & Red, Green, Blue & Rojo, verde, azul \\
RHS & Right Hand Side & Lado derecho (de la ecuación) \\
RNA & Ribonucleic Acid & Ácido ribonucleico \\
RoPE & Rotary Positional Embeddings  & Embeddings de posición rotacionales \\
SAE & Sparse Autoencoder  & Autocodificador esparso \\
SDF & Signed Distance Field & Campo de distancia con signo \\
SFM & Structure From Motion & Estructura desde el movimiento \\
SGD & Stochastic Gradient Descent & Descenso de gradiente estocástico \\
SiLU & Sigmoid Linear Unit & Unidad lineal sigmoide \\
SITA & Shape Image Text Aligned & Alineado Forma Imagen Texto \\
SLAM & Simultaneous Localization And Mapping & Localización y mapeo simultáneos \\
SMPL & Skinned Multi-Person Linear & Modelo lineal multi-persona con piel \\
SNR & Signal to Noise Ratio & Relación señal-ruido \\
SSA & Subspace Self-Attention & Autoatención por subespacios \\
STFT & Short-Time Fourier Transform & Transformada de Fourier de tiempo reducido \\
SURE & Stein's Unbiased Risk Estimator & Estimador de riesgo no sesgado de Stein \\
SVD & Singular Value Decomposition & Descomposición en valores singulares \\
SVM & Support Vector Machine & Máquina de vectores de soporte \\
TDS & Twisted Diffusion Sampler & Muestreo torcido de difusión \\
ToST & Token Statistics Transformer & Transformer de estadística de tokens \\
TSSA & Token Statistics Self-Attention & Autoatención de estadística de tokens \\
TTT & Test Time Training & Entrenamiento en tiempo de prueba \\
TV & Total Variation & Variación total \\
ULIP & unified language, image, and point cloud learning & Aprendizaje unificado de lenguaje, imagen y nubes de puntos \\
VAE & variational autoencoder & Autocodificador variacional \\
VE & Variance-Exploding (process) & Proceso de difusión de varianza explosiva \\
ViT & Vision Transformer & Transformador de visión \\
VP & Variance-Preserving (process) & Proceso de difusión de varianza constante \\
VQ & Vector Quantization & Cuantización vectorial \\
WLOG & without loss of generality & Sin pérdida de generalidad \\
\hline
%\end{tabular}
\caption{Acrónimos usados en este texto.}
\label{tab:basic}
%\end{table}
\end{longtable}

\end{document}
